% ---------------------------------------------------------------
% Notation for the criterion layer (subgroup reflection, peripheral
% isolation, tile/base reduction).  Assumes a standard amsthm preamble:
%   \newtheorem{theorem}{Theorem}[section]
%   \newtheorem{lemma}[theorem]{Lemma}
%   \newtheorem{proposition}[theorem]{Proposition}
%   \newtheorem{corollary}[theorem]{Corollary}
%   \theoremstyle{definition}\newtheorem{definition}[theorem]{Definition}
%   \theoremstyle{remark}\newtheorem{remark}[theorem]{Remark}
% ---------------------------------------------------------------
\subsection*{Notation and standing conventions}

\begin{itemize}
\item \(K\Subset U\) means \(K\) is compactly contained in \(U\).  A
      \emph{full continuum} is a nonempty compact connected set
      \(K\subset\mathbb C\) with connected complement
      \(\mathbb C\setminus K\).
\item \(f:U^1\to U^0\) is an ordinary quadratic-like restriction of a
      quadratic polynomial; \(U^r=f^{-r}(U^0)\), so
      \(U^{r+1}\subset U^r\) and each \(f:U^{r+1}\to U^r\) is
      quadratic-like (Kahn~I, \S6, p.~31).  For a
      primitively renormalizable stage of period \(p>1\) we write
      \(F=f^p:U^p\to U^0\), proper of degree \(2^p\), and
      \(i^p:U^p\hookrightarrow U^0\) for the physical inclusion.
\item \(K_p\) is the little filled Julia set, \(K_p(j)=f^j(K_p)\) for
      \(0\le j<p\), and \(\mathcal K_p=\bigcup_{j=0}^{p-1}K_p(j)\);
      primitivity means the \(K_p(j)\) are pairwise disjoint (Kahn~I,
      \S1.1).  We abbreviate \(A=K_p\),
      \(B=\mathcal K_p\).
\item The canonical-localization marked surfaces are
      \(T'=(U^p\setminus F^{-1}B)\cup A\) and
      \(T=(U^0\setminus B)\cup A\), with complementary surfaces
      \(\Omega_s=U^p\setminus F^{-1}B\) and
      \(\Omega_t=U^0\setminus B\)
      (Kahn~I, \S\S7.4.2--7.4.3, pp.~45--46).
\item For a subgroup \(H\le\pi_1(X,x)\) of a connected, locally
      path-connected, semilocally simply connected space, \(X_H\to X\)
      denotes the connected based cover corresponding to \(H\).  The
      \emph{localization subgroups}
      \(H_s\le\pi_1(\Omega_s,z_*)\), \(H_t\le\pi_1(\Omega_t,z_*)\) are
      the images of the fundamental groups of deleted local
      neighborhoods of the marked continuum \(A\)
      (Kahn~I, \S7.4.1, Theorem~7.12, p.~44).  Kahn requires of a local
      neighborhood only that every component of its complement of \(A\)
      be an annulus; the specialization to \emph{Jordan} neighborhoods
      made in this paper is our own, and is justified by
      Lemma~\ref{lem:local-pair}.
\item Boldface domains \(\mathbf U^m\) are the stages of the exact
      equalizer tower of Section~\ref{sec:forced} over the canonical
      seed, with complementary annuli
      \(A_m=\mathbf U^m\setminus K^m\).  Boldface domains
      \(\mathbf U_{\mathrm{KL}}^m\), with
      \(A^{\mathrm{KL}}_m=\mathbf U_{\mathrm{KL}}^m\setminus K(f)\),
      are the \(\psi\)-quadratic-like (bold) domains on which the
      Kahn--Lyubich pseudo-puzzle of KL~II, \S3.2 is run.  These two
      families of symbols are kept strictly distinct: comparing them is
      the entire content of the strong stage-alignment hypothesis
      (Definition~\ref{def:strong-alignment}), and they may be
      identified only where that hypothesis has been declared to be in
      force.  \(\boldsymbol\iota\) denotes
      the inclusion-induced pseudo-piece \emph{immersion} and
      \(\mathbf g\) the
      dynamical branched-covering arrow; these are distinct arrows
      (KL~II, Lemma~3.4, pp.~12--13).  We follow KL in reserving the
      word ``inclusion'' for the puzzle-piece relation \(Y\subset Z\):
      the ambient arrow is an immersion, and is never proved to be an
      embedding.
\item \((q,\nu)\) are what KL~II call the \emph{decoration parameters}
      (KL~II, Corollary~5.14 and Theorem~5.15, pp.~23--24), \(q\) being
      the denominator of the combinatorial rotation number at the
      \(\alpha\)-fixed point (KL~II, \S2.1, p.~5).  Where this paper
      says ``decoration period \(q\) and depth index \(\nu\)'' it means
      exactly KL's decoration parameters.
\item For a connected annulus, \(\pi_\bullet:\widetilde A_\bullet\to
      A_\bullet\) is a universal cover with positively oriented deck
      generator \(T_\bullet\); the (annular) \emph{degree} of a map of
      annuli is the integer it induces on \(\pi_1\cong\mathbb Z\) with
      respect to the chosen orientations.
\item \(\operatorname{Ind}_x(\gamma)\) is the winding number of a loop
      \(\gamma\) about a point \(x\in\mathbb C\setminus|\gamma|\).
\item Citations of the form ``Kahn~I, \S X / Theorem~N, p.~P'' refer to
      \cite{KahnPrimitive06} in the numbering of arXiv:math/0609045v3
      (28 December 2023); citations of the form ``KL~II, Lemma~N,
      p.~P'' refer to \cite{KLDecorations08} in the numbering of
      Ann.\ Sci.\ \'Ec.\ Norm.\ Sup\'er.\ (4) \textbf{41} (2008)
      57--84 and of the identically sectioned arXiv:math/0609046v2.
      The SHA-256--locked local \TeX{} archives and their line-number
      concordance are recorded in Appendix~\ref{app:sourcelock} only.
\end{itemize}